\section{Risk Aversion Endogen: Pendekatan Sigmoid}\label{appendix-b}

\hypertarget{urgensi-mengapa-sigmoid}{%
\subsection{Urgensi: Mengapa Sigmoid?}\label{urgensi-mengapa-sigmoid}}

Dalam model Markowitz statis, \(\lambda\) (bobot return) seringkali dipilih secara subjektif (misal \(\lambda=0.5\)). Namun, pasar yang riil bersifat dinamis. Kita membutuhkan fungsi yang bisa mengubah rasio performa aset (\(\mu/\sigma\)) menjadi bobot keputusan yang halus namun memiliki batas yang jelas. Fungsi Sigmoid adalah kandidat sempurna karena ia memetakan input tak terbatas menjadi output dalam rentang \((0, 1)\).

\begin{center}\rule{0.5\linewidth}{0.5pt}\end{center}

\hypertarget{aksioma-intuisi-psikologi-keputusan}{%
\subsection{Aksioma \& Intuisi: Psikologi Keputusan}\label{aksioma-intuisi-psikologi-keputusan}}

Ada tiga aksioma dasar yang melandasi penggunaan fungsi ini: 
\begin{enumerate}
  \item \textbf{Aksioma Signal-to-Noise:} Investor tidak melihat keuntungan (\(\mu\)) secara terisolasi, melainkan selalu membandingkannya dengan derau/risiko (\(\sigma\)). Rasio \(\mu/\sigma\) adalah penggerak utama keputusan.
  \item \textbf{Aksioma Saturasi:} Nafsu terhadap keuntungan tidak tumbuh linear tanpa batas. Ada titik di mana tambahan keuntungan tidak lagi menambah ``nafsu'' secara signifikan (titik jenuh/saturasi).
  \item   \textbf{Aksioma Batas:} Parameter \(\lambda\) harus berada di antara 0 (fokus total pada risiko) dan 1 (fokus total pada return).
\end{enumerate}

\begin{center}\rule{0.5\linewidth}{0.5pt}\end{center}

\hypertarget{reduksionisme-kasus-minimal-mu-sigma}{%
\subsection{\texorpdfstring{Reduksionisme: Kasus Minimal (\(\mu = \sigma\))}{Reduksionisme: Kasus Minimal (\textbackslash mu = \textbackslash sigma)}}\label{reduksionisme-kasus-minimal-mu-sigma}}

Mari kita lihat apa yang terjadi jika risiko suatu aset tepat sama dengan keuntungan yang diharapkan.

\hypertarget{a.-kondisi-seimbang}{%
\subsubsection{Kondisi Seimbang}\label{a.-kondisi-seimbang}}

Jika \(\mu = \sigma\), maka rasio \(\mu/\sigma = 1\). \[ \lambda = \frac{1}{1 + e^{-1}} \qquad (1) \]

\begin{quote}
\textbf{Visualisasi (1): Jembatan Logika} Secara numerik: \[ \lambda \approx \frac{1}{1 + 0.367} \approx 0.731 \] Karena \(\mu/\sigma\) positif dan cukup kuat, sistem ``condong'' untuk mengejar return (\(\lambda > 0.5\)). Sebaliknya, jika \(\mu = 0\) (tidak ada return), maka \(\lambda = 0.5\) (netral).
\end{quote}

\begin{center}\rule{0.5\linewidth}{0.5pt}\end{center}

\hypertarget{jembatan-logika-formalisme-fungsi-aktivasi}{%
\subsection{Jembatan Logika: Formalisme Fungsi Aktivasi}\label{jembatan-logika-formalisme-fungsi-aktivasi}}

Dalam \emph{Machine Learning}, sigmoid digunakan sebagai gerbang (\emph{gate}). Dalam VQE Portofolio, \(\lambda\) adalah gerbang yang menentukan seberapa besar ``angin return'' boleh meniup partikel kita keluar dari ``lembah risiko''.

\hypertarget{a.-persamaan-umum}{%
\subsubsection{Persamaan Umum}\label{a.-persamaan-umum}}

\[ \lambda(\mu, \sigma) = \frac{1}{1 + e^{-(\mu/\sigma)}} \qquad (2) \]

\begin{center}\rule{0.5\linewidth}{0.5pt}\end{center}

\hypertarget{derivasi-indeks-persamaan}{%
\subsection{Derivasi \& Indeks Persamaan}\label{derivasi-indeks-persamaan}}

Mari kita turunkan bagaimana perubahan kecil pada risiko (\(\Delta \sigma\)) memengaruhi selera investasi kita.

\hypertarget{a.-sensitivitas-selera-terhadap-risiko}{%
\subsubsection{Sensitivitas Selera terhadap Risiko}\label{a.-sensitivitas-selera-terhadap-risiko}}

Turunan \(\lambda\) terhadap \(\sigma\) menunjukkan seberapa cepat selera kita berubah saat pasar mulai bergejolak: \[ \frac{\partial \lambda}{\partial \sigma} = \frac{e^{-(\mu/\sigma)} \cdot (-\mu/\sigma^2)}{(1 + e^{-(\mu/\sigma)})^2} \qquad (3) \]

\begin{quote}
\textbf{Visualisasi (3): Perhitungan Linear} Kita sederhanakan menggunakan properti turunan sigmoid \(f'(x) = f(x)(1-f(x))\): \[ \frac{\partial \lambda}{\partial \sigma} = \lambda(1-\lambda) \cdot \left(\frac{\mu}{\sigma^2}\right) \] Misalkan pada titik netral \(\lambda = 0.5\) dan \(\mu/\sigma^2 = 2\): \[ \frac{\partial \lambda}{\partial \sigma} = (0.5)(0.5) \cdot 2 = 0.5 \] Hal ini menggambarkan bahwa pada kondisi netral, perubahan volatilitas memiliki dampak 50\% terhadap pergeseran selera investasi kita.
\end{quote}

\begin{center}\rule{0.5\linewidth}{0.5pt}\end{center}

\hypertarget{verifikasi-parameter-batas-ekstrem}{%
\subsection{Verifikasi \& Parameter: Batas Ekmenstrem}\label{verifikasi-parameter-batas-ekstrem}}

Kita uji model ini pada tiga kondisi pasar yang berbeda:

\begin{longtable}[]{@{}
  >{\raggedright\arraybackslash}p{(\columnwidth - 6\tabcolsep) * \real{0.2500}}
  >{\raggedright\arraybackslash}p{(\columnwidth - 6\tabcolsep) * \real{0.2500}}
  >{\raggedright\arraybackslash}p{(\columnwidth - 6\tabcolsep) * \real{0.2500}}
  >{\raggedright\arraybackslash}p{(\columnwidth - 6\tabcolsep) * \real{0.2500}}@{}}
\toprule\noalign{}
\begin{minipage}[b]{\linewidth}\raggedright
Kondisi Pasar
\end{minipage} & \begin{minipage}[b]{\linewidth}\raggedright
Rasio \(\mu/\sigma\)
\end{minipage} & \begin{minipage}[b]{\linewidth}\raggedright
Nilai \(\lambda\)
\end{minipage} & \begin{minipage}[b]{\linewidth}\raggedright
Perilaku VQE
\end{minipage} \\
\midrule\noalign{}
\endhead
\bottomrule\noalign{}
\endlastfoot
\textbf{Pasar Euphoria} & \(\mu \gg \sigma\) & \(\lambda \to 1\) & Memilih aset return tertinggi, abaikan risiko. \\
\textbf{Pasar Netral} & \(\mu = 0\) & \(\lambda = 0.5\) & Keseimbangan antara keamanan dan profit. \\
\textbf{Pasar Panik} & \(\mu \ll 0\) & \(\lambda \to 0\) & Memilih aset paling aman (varians terendah). \\
\end{longtable}

\begin{quote}
\textbf{Visualisasi (6): Contoh Perhitungan} Jika pasar sedang jatuh (\(\mu = -10, \sigma = 2\)): \[ \text{Rasio} = -5 \implies \lambda = \frac{1}{1 + e^5} \approx \frac{1}{1 + 148} \approx 0.006 \] Nilai \(\lambda \approx 0\) ini akan membuat Hamiltonian VQE didominasi oleh suku risiko \(x^T \Sigma x\).
\end{quote}

\begin{center}\rule{0.5\linewidth}{0.5pt}\end{center}
